\section*{补证：Hurwitz判据为何成立}
\addcontentsline{toc}{section}{补证：Hurwitz判据为何成立}
\label{sec:hurwitz-proof}

\noindent\textbf{编者补证。}
原书§9只陈述定理7，另引文献给出证明。
为使本稿可独立使用，这里补上证明；它不冒充原书已有的段落。
证明采用与控制理论相通的Routh降阶思路，再把降阶条件写成Hurwitz顺序主子式。
沿用本章已使用的代数事实：固定次数多项式的根作为一个带重数的集合，随系数连续变化。
这不要求能把所有根在全参数范围内逐个编号为单值函数。

\subsection*{1. 奇偶部分与降阶多项式}

设
\[
 P(s)=a_0s^n+a_1s^{n-1}+\cdots+a_n,\qquad a_0>0,
\]
各$a_j$均为实数，且$n\ge2$。
先假定$a_1>0$。把多项式按幂次的奇偶性分成
\[
 U(s)=a_0s^n+a_2s^{n-2}+a_4s^{n-4}+\cdots,
 \qquad
 V(s)=a_1s^{n-1}+a_3s^{n-3}+a_5s^{n-5}+\cdots.
\]
令$c=a_0/a_1>0$，定义
\[
 W(s)=U(s)-csV(s),\qquad Q(s)=V(s)+W(s).
\]
$U$与$csV$的最高次项相消，所以$\deg W\le n-2$，
而$Q$恰为$n-1$次多项式，首项系数$a_1>0$。
写开即为
\[
 Q(s)=a_1s^{n-1}
 +\left(a_2-\frac{a_0a_3}{a_1}\right)s^{n-2}
 +a_3s^{n-3}
 +\left(a_4-\frac{a_0a_5}{a_1}\right)s^{n-4}
 +\cdots.
\]
系数越界时取零。它就是Routh过程的一次降阶。

\subsection*{2. 降阶为什么保持左半平面稳定性}

我们证明
\begin{equation}
 P\text{稳定}
 \quad\Longleftrightarrow\quad
 a_1>0\ \text{且}\ Q\text{稳定}.
 \label{eq:routh-reduction}
\end{equation}
先注意：若$P$稳定，实根给出因子$s+r$，$r>0$；
一对共轭根$-\alpha\pm i\beta$给出
$s^2+2\alpha s+\alpha^2+\beta^2$，$\alpha>0$。
这些因子的系数均为正，故乘积的系数均为正，特别有$a_1>0$。

在$a_1>0$的条件下，考虑实参数族
\[
 P_\tau(s)=W(s)+V(s)+\tau sV(s),\qquad \tau>0.
\]
其中$P_c=P$，而当$\tau\downarrow0$时$P_\tau\to Q$。
对实数$\omega$，$W(i\omega)$与$i\omega V(i\omega)$
都属于同一个实直线$i^n\R$，
$V(i\omega)$则属于与它垂直的实直线$i^{n-1}\R$。
因此
\[
 P_\tau(i\omega)=0
 \quad\Longleftrightarrow\quad
 V(i\omega)=0\ \text{且}\ W(i\omega)=0
 \quad\Longleftrightarrow\quad Q(i\omega)=0.
\]
这个条件与$\tau$无关。
所以，只要$P$或$Q$中一个没有虚轴根，整个$\tau>0$的多项式族都没有虚轴根。

在任意闭参数区间$0<\tau_1\le\tau\le\tau_2$上，
$P_\tau$的首项系数为$\tau a_1$，不会消失。
根随系数连续变化，并且在这个闭参数区间内有统一的根模上界。
因此右半平面的根数若发生改变，必有某个根经过虚轴。
既然虚轴根不存在，右半平面根数就在所有$\tau>0$上保持不变。

还需要检查$\tau\downarrow0$时那个因降阶而“走向无穷”的根。
$Q$有$n-1$个有限根，按重数计；$P_\tau$有$n-1$个根趋向这些根，
余下一个根记为$r_\tau$。
由根与系数的关系，$P_\tau$全部根之和恰为$-1/\tau$。
有限的$n-1$个根之和有界，所以
\[
 r_\tau=-\frac1\tau+O(1).
\]
由于系数为实数，非实根必须成共轭对；
只有一个根离开所有有界区域，因此这个根在充分小的$\tau$下必为实数，
并且为负。
关于有限根的收敛，也可把
$P_\tau(s)=\tau a_1\prod_{j=1}^n(s-r_j(\tau))$
局限在包含$Q$全部根的一个大圆内来理解：圆内的根按重数保持为$n-1$个，
剩下的一个根离开圆。

若$Q$稳定，则充分小$\tau>0$时，趋向$Q$根的$n-1$个根都在左半平面，
额外的实根也为负。因此小$\tau$时$P_\tau$稳定；
再由没有虚轴穿越，得到$P=P_c$稳定。

反过来，若$P$稳定，则所有$\tau>0$的$P_\tau$都稳定。
令$\tau\downarrow0$，$Q$的有限根只能位于闭左半平面。
但若$Q$有虚轴根，上面的等价式会使$P$也有同一个虚轴根，矛盾。
因此$Q$所有根实部严格为负。降阶等价式得证。

\subsection*{3. Hurwitz顺序主子式的递推}

定义
\[
 \mat{H}(P)=(a_{2j-i})_{i,j=1}^n
 =
 \begin{pmatrix}
 a_1&a_3&a_5&\cdots\\
 a_0&a_2&a_4&\cdots\\
 0&a_1&a_3&\cdots\\
 0&a_0&a_2&\cdots\\
 \vdots&\vdots&\vdots&\ddots
 \end{pmatrix},
\]
并令$\Delta_k(P)$为其左上角$k\times k$子矩阵的行列式。
对任意$k\ge2$，在这个子矩阵中，
把每个偶数行减去其上一个奇数行的$c=a_0/a_1$倍。
这不改变行列式，并把第一列变成$(a_1,0,\ldots,0)^{\mathsf T}$。
按第一列展开，剩余的子矩阵恰为$\mat{H}(Q)$的左上角$(k-1)\times(k-1)$部分。
故
\begin{equation}
 \Delta_1(P)=a_1,\qquad
 \Delta_k(P)=a_1\Delta_{k-1}(Q)\quad(2\le k\le n).
 \label{eq:hurwitz-recursion}
\end{equation}
例如第二阶给出
\[
 \Delta_2(P)
 =a_1\left(a_2-\frac{a_0a_3}{a_1}\right)
 =a_1a_2-a_0a_3,
\]
正是降阶后首两项系数所决定的第一个条件。

\subsection*{4. 归纳完成定理7}

当$n=1$时，$P(s)=a_0s+a_1$，其唯一根为$-a_1/a_0$；
在$a_0>0$下，稳定当且仅当$a_1>0$，也就是$\Delta_1(P)>0$。

假设结论对$n-1$次多项式成立。
若$n$次多项式$P$稳定，由\eqref{eq:routh-reduction}有$a_1>0$且$Q$稳定。
归纳假设给出$\Delta_j(Q)>0$，$1\le j\le n-1$。
再由\eqref{eq:hurwitz-recursion}得到$\Delta_k(P)>0$，$1\le k\le n$。

反过来，若所有$\Delta_k(P)>0$，则$a_1=\Delta_1(P)>0$，
所以降阶多项式$Q$可定义，并由\eqref{eq:hurwitz-recursion}得到
$\Delta_j(Q)>0$，$1\le j\le n-1$。
归纳假设说明$Q$稳定，降阶等价式再说明$P$稳定。
于是
\[
 \boxed{\quad P\text{的全部根满足}\ \Re s<0
 \quad\Longleftrightarrow\quad
 \Delta_1(P),\ldots,\Delta_n(P)>0.\quad}
\]
证明完成。

这里判定的是严格左半平面稳定性。
某个顺序主子式为零时，不能把上述严格不等式改成非严格不等式继续套用；
边界根、零行等情形需要另行分析。
